\section{A uniform bound for transient reflection classes}\label{sec:transients}

We now bound the total dimension of reflection-equivariant classes at a
finite cut modulo those coming from the full loop space.  Ordinary death
forces such a class, away from finitely many index bands, to propagate
through a long string of powers of $w$.  The finite-cut budget rules out
that string.  Bounds in the remaining degrees then give a bound on the
total dimension.

Throughout this section $n\ge3$, $d=n-1$, and we assume that there are
only finitely many prime closed-geodesic images.  Let $\ell_i>0$ and
$\Delta_i\ge0$ be their prime lengths and mean indices, as in
Theorem~\ref{thm:local-control}; thus
\begin{equation}\label{trans:index-slope}
 |I_{i,m}-m\Delta_i|\le n,
 \qquad \beta_i=\Delta_i/\ell_i.
\end{equation}
For every positive regular cut $a$, retain
\[
 \begin{gathered}
 B_a^j=H_R^j(X_a,M),\qquad B_\infty^j=H_R^j(\Lambda,M),\\
 Q_a^j=\coker(B_\infty^j\longrightarrow B_a^j),
 \qquad T(a)=\sum_{j\ge0}\dim Q_a^j.
 \end{gathered}
\]
All coefficients are $\kk$.  The local crossing bounds imply that
$\bigoplus_j B_a^j$ is finite-dimensional at each such cut.

\begin{theorem}[Bounded total transients]\label{thm:bounded-transients}
Under the finite-prime-image assumption, there is a constant $C$,
independent of the positive regular cut $a$, such that
\begin{equation}\label{trans:total-bound}
 T(a)\le C.
\end{equation}
\end{theorem}

We use four previously established inputs.  Theorem~\ref{thm:ordinary-death}
provides a single filling increment $D$ in ordinary homology relative
to $M$, independent of both degree and cut.
Theorem~\ref{thm:local-control} bounds local ranks and their index
supports.  Theorem~\ref{thm:global-generators} gives
$w^nB_\infty=0$.  Finally, Theorem~\ref{thm:finite-cut-budget}
controls the discrepancy between actual ambient degree and root degree.
The ordinary rank and resonance statements used below are those of
Proposition~\ref{prop:sphere-powers}; write $\alpha>0$ for its
resonance slope.  Constants in this section may depend on the metric,
$n$, and the finite list of prime images, but never on $a$ or $j$.

\subsection{Degree gaps and multiplication by \texorpdfstring{$w$}{w}}

Let
\[
 \mathcal S=\{0,\alpha,\beta_i:\text{prime labels }i\}
           =\{s_0<s_1<\cdots<s_N\}
\]
be the list of distinct slopes.  Thus $s_0=0$ and $N\ge1$.
For a fixed margin $c>0$ and consecutive slopes $s<s'$, define the
integer degree gap
\[
 \mathcal I_{s,s'}(a)
 =\{j\in\mathbb Z_{\ge0}:sa+c<j<s'a-c\}.
\]
Since $\alpha$ is among the slopes, each gap lies entirely below or
entirely above $\alpha a$.  We call these gaps present and absent,
respectively.

\begin{lemma}[Comparison on degree gaps]\label{trans:degree-gaps}
There is a fixed margin $c$ such that, for every positive regular $a$
and every consecutive pair $s<s'$ in $\mathcal S$, the following
statements hold.
\begin{enumerate}
 \item In a present gap, $B_\infty^j\to B_a^j$ is injective for every
       $j\in\mathcal I_{s,s'}(a)$.
 \item In an absent gap, $B_\infty^j\to B_a^j$ is zero for every
       $j\in\mathcal I_{s,s'}(a)$.
 \item Whenever $j,j+1\in\mathcal I_{s,s'}(a)$, multiplication by $w$
       induces an isomorphism $Q_a^j\to Q_a^{j+1}$.
\end{enumerate}
Moreover $B_a^j=0$ for $j>s_Na+c$.
\end{lemma}

\begin{proof}
We first prove the ordinary comparison away from fixed-width bands.
Choose a regular endpoint $a_+\in(a+D,a+D+1)$.
A local group at an event $m\ell_i\in[a,a_+]$ can contribute in
degree $j$ or its neighboring degrees only if
\[
 |j-I_{i,m}|\le C_{\rm loc}.
\]
By \eqref{trans:index-slope}, this implies
\begin{equation}\label{trans:horizon-band}
 |j-\beta_i a|
 \le C_{\rm loc}+n+\beta_i(D+1).
\end{equation}
The right-hand side is bounded uniformly over the finite prime list.
Thus, outside fixed-width bands about the $\beta_i a$, the ordinary
map
\[
 H_j(X_a,M)\longrightarrow H_j(X_{a_+},M)
\]
is injective, by the crossing exact sequences.  A class in its source
that is zero globally is already zero in $X_{a+D}$ by
Theorem~\ref{thm:ordinary-death}.  Injectivity therefore makes the
map $H_j(X_a,M)\to H_j(\Lambda,M)$ injective.  Field duality gives
surjectivity of ordinary cohomology restriction.

In degrees above a fixed low threshold, the ordinary groups relative
to $M$ identify with the absolute groups, by the evaluation retraction
and the cohomology of the constant sphere.  In these degrees
Proposition~\ref{prop:sphere-powers} gives global rank at most one.
Its resonance estimate implies that every nonzero global class
restricts nontrivially to $X_a$ when $j<\alpha a-c_0$, and restricts
to zero when $j>\alpha a+c_0$, for a fixed $c_0$.
Together with the preceding surjectivity, this proves
\begin{equation}\label{trans:ordinary-gaps}
 \begin{cases}
 H^j(\Lambda,M)\longrightarrow H^j(X_a,M)
       \text{ is an isomorphism},&j<\alpha a-c_0,\\
 H^j(X_a,M)=0,&j>\alpha a+c_0,
 \end{cases}
\end{equation}
provided $j$ also lies outside the fixed low-degree interval and the
bands in \eqref{trans:horizon-band}.  When the global group is zero,
surjectivity gives the first assertion as an isomorphism between zero
groups.  Increase $c_0$ to absorb all these restrictions.

Consider a present gap and put
$P_a^j=H_R^j(\Lambda,X_a)$ for this proof.
Ordinary triple exactness and \eqref{trans:ordinary-gaps} give
$H^j(\Lambda,X_a)=0$ whenever the needed two neighboring ordinary
restrictions are isomorphisms.  The sign-line Gysin sequence of
Proposition~\ref{prop:gysin} consequently makes
$w:P_a^j\to P_a^{j+1}$ an isomorphism in the corresponding interior.
After allowing $n$ further degrees below, every element of $P_a^j$
is divisible by $w^n$ there.  Its image in $B_\infty^j$ is zero,
since $w^nB_\infty=0$.  Applying this at $j$ and $j+1$ to the triple
$(\Lambda,X_a,M)$ gives a natural short exact sequence
\begin{equation}\label{trans:present-sequence}
 0\longrightarrow B_\infty^j\longrightarrow B_a^j
   \longrightarrow P_a^{j+1}\longrightarrow0.
\end{equation}
The quotient identification commutes with $w$.  It gives both the
claimed injection and the isomorphisms on $Q_a$ in a present gap.

In an absent gap, \eqref{trans:ordinary-gaps} and sign Gysin make
$w:B_a^j\to B_a^{j+1}$ an isomorphism in the interior.  Allowing $n$
further degrees above makes $w^n$ injective there.  The image of
$B_\infty^j$ is killed by that power, so it is zero.  Hence $Q_a^j=B_a^j$
and the claimed isomorphisms also hold in an absent gap.

Choose the margin $c$ larger than $c_0+2n+6$ and the fixed low-degree
threshold.  Then all neighboring degrees and all powers just used lie
in the required untrimmed ordinary ranges whenever the statements of
the lemma apply.  This is one fixed choice for all cuts and gaps.
Finally, any reflection local contribution below cut $a$ has
\[
 j\le I_{i,m}+C_{\rm loc}
   \le \beta_i a+n+C_{\rm loc}\qquad(m\ell_i<a).
\]
The relative groups start at zero below the first event.  The local
crossing sequences therefore give $B_a^j=0$ above $s_Na+c$, after
one final increase of $c$.
\end{proof}

\subsection{The finite-cut budget excludes long strings}

\begin{lemma}[No transients inside a long degree gap]\label{trans:no-interior}
For all sufficiently large regular cuts $a$ with $2a$ regular,
$Q_a^j=0$ on every gap $\mathcal I_{s,s'}(a)$ of
Lemma~\ref{trans:degree-gaps}.
\end{lemma}

\begin{proof}
Fix consecutive slopes $s<s'$ and write
\[
 A=sa,\qquad r=(s'-s)a.
\]
There is a fixed $\delta>0$ such that $r\ge\delta a$ for every such
gap, since the list of distinct slopes is finite.  Suppose that some
$Q_a^j$ in this gap is nonzero.  The isomorphisms in
Lemma~\ref{trans:degree-gaps} make it nonzero throughout the integer
gap.  For large $a$, choose
\[
 l_0=\lceil A+0.10r\rceil
\]
and a class $y_{l_0}\in B_a^{l_0}$ with nonzero image in $Q_a^{l_0}$.
Set $y_l=w^{l-l_0}y_{l_0}$ through
$l\le\lfloor A+0.90r\rfloor$.  Every resulting quotient class is
nonzero.  In particular, with
\[
 \mathcal L=\mathbb Z\cap[A+0.15r,A+0.20r],
 \qquad q=\lfloor0.50r\rfloor,
\]
the classes $y_l$, $l\in\mathcal L$, are independent in
$\bigoplus_j B_a^j$, as are their powers $w^q y_l$.
Indeed they occupy distinct root degrees, and their quotient classes
remain nonzero between degrees $A+0.65r-O(1)$ and $A+0.70r$.
Figure~\ref{fig:degree-gap} displays the intervals used below.

At the ambient cut $2a$, let $E_J$ and $P_k$ be the actual normalized
restriction image and root-degree projection from
Theorem~\ref{thm:finite-cut-budget}.  Recall its nonnegative terms
\[
 D_{k,J}=\dim P_k\Bigl(\bigoplus_l B_a^l\Bigr)-\dim P_kE_J,
 \qquad H_{k,J}=\dim(E_J\cap\ker P_k).
\]
In particular that theorem gives
\begin{equation}\label{trans:used-budget}
 \sum_{k\ge1}\bigl(D_{k,2k}+H_{k,2k}\bigr)\le C(1+a).
\end{equation}
Average $D_{k,2k}$ over the integers
$k\in[A+0.30r,A+0.35r]$.  There are a fixed positive multiple of $r$
such integers for large $r$, so some $k$ satisfies
$D_{k,2k}\le C_1$, independently of $a$, because $r\ge\delta a$.
The space
\[
 W_0=\operatorname{span}\{y_l:l\in\mathcal L\}
\]
lies in the root degrees at most $k$.  Thus
\[
 \dim(W_0\cap P_kE_{2k})
 \ge\dim W_0-D_{k,2k}\ge c_1r
\]
for a fixed $c_1>0$ and sufficiently large $r$.

Choose a basis $z_1,\ldots,z_\nu$ of this intersection and actual
ambient lifts
\[
 u_\mu\in H_K^{2k}(X_{2a},M).
\]
The normalized fixed restriction of $u_\mu$ has the form
$z_\mu+c_\mu$, where every root label in $c_\mu$ is greater than $k$.
Multiply the actual classes by $w^q$.  The selected terms
$w^q z_\mu$ remain independent, since $w^q$ is injective on $W_0$.
Their labels are at most $A+0.70r$, whereas the other terms have
labels greater than
\[
 k+q\ge A+0.80r-1
\]
or vanish.  The latter terms cannot cancel any of the selected ones.
Moreover all labels of these normalized images exceed $A+0.64r$
for sufficiently large $r$.

For any integer $k'\in[A+0.61r,A+0.64r]$, there is enough actual
degree to multiply further by a nonnegative power of $v$, since
\begin{equation}\label{trans:degree-room}
 2k+q\le2A+1.20r<2A+1.22r\le2k'.
\end{equation}
\begin{figure}[tbp]
\centering
\begin{tikzpicture}[x=10.8cm,y=.72cm,>=Stealth,
  every node/.style={font=\small},
  selected/.style={blue!65!black,line width=4pt},
  cutoff/.style={teal!70!black,line width=4pt},
  remainder/.style={gray,line width=3pt}]
  \draw[gray!55,->] (0,-.65) -- (1.025,-.65);
  \foreach \x/\lab in {0/0,.25/{.25},.5/{.50},.75/{.75},1/1} {
    \draw[gray!55] (\x,-.72) -- (\x,-.58);
    \node[below=3pt] at (\x,-.72) {$\lab$};
  }
  \node[below=20pt] at (.5,-.72) {normalized root degree $(l-A)/r$};

  \node[left=7pt,align=right] at (0,3.6) {selected classes\\in $W_0$};
  \draw[selected] (.15,3.6) -- (.20,3.6);
  \node[above=4pt] at (.175,3.6) {$[.15,.20]$};

  \node[left=7pt,align=right] at (0,2.5) {first cutoff\\$k$};
  \draw[cutoff] (.30,2.5) -- (.35,2.5);
  \node[above=4pt] at (.325,2.5) {$[.30,.35]$};

  \node[left=7pt,align=right] at (0,1.4) {later cutoffs\\$k'$};
  \draw[cutoff] (.61,1.4) -- (.64,1.4);
  \node[above=4pt] at (.625,1.4) {$[.61,.64]$};

  \node[left=7pt,align=right] at (0,.3) {lifted images\\after $w^q$};
  \draw[selected] (.65,.3) -- (.70,.3);
  \node[above=4pt] at (.675,.3) {$[.65,.70]$};
  \draw[remainder,->] (.805,.3) -- (.98,.3);
  \node[above=4pt,text=gray!80!black] at (.895,.3) {possible tail $>.80$};
\end{tikzpicture}
\caption{The degree intervals in Lemma~\ref{trans:no-interior}, normalized
by the gap width $r$; the drawing suppresses rounding errors of order
$r^{-1}$. Multiplication by $w^q$, with $q/r\approx .50$, moves the
selected root degrees above every later cutoff $k'$ while keeping them
separate from the other terms. Multiplication by $v$ then raises actual
degree without changing root labels; the inequality
$2k+q<2k'$ provides the required degree room.}
\label{fig:degree-gap}
\end{figure}

The classes
\[
 v^{\,2k'-(2k+q)}w^q u_\mu
       \in H_K^{2k'}(X_{2a},M)
\]
have unchanged normalized fixed images: multiplication by $v$ changes
the actual degree and leaves the root coefficients unchanged.
Multiplication by $w$ raises each root label by one, as used above.
Thus the images are independent and lie entirely above root cutoff
$k'$, proving
\[
 H_{k',2k'}\ge c_1r
 \quad\text{for every integer }k'\in[A+0.61r,A+0.64r].
\]
There are at least $c_2r$ such cutoffs for large $r$.  The left side of
\eqref{trans:used-budget} is therefore at least
$c_1c_2r^2\ge c_1c_2\delta^2a^2$, contradicting its linear upper bound.
This proves the assertion for the chosen gap.  The finite number of
gaps allows one common lower threshold for $a$.
\end{proof}

\subsection{Bounding the remaining degrees}

\begin{proof}[Proof of Theorem~\ref{thm:bounded-transients}]
First let $a$ be large with $2a$ regular.  By
Lemma~\ref{trans:no-interior}, transient classes can occur only within
distance $c$ of one of the finitely many degrees $s_i a$.
Lemma~\ref{trans:degree-gaps} excludes degrees above the last such
band.  Thus only a bounded number of integer degrees can contribute.
We now bound the dimension in each of them uniformly.

For high degrees, the reflection local ranks of
Theorem~\ref{thm:local-control} give
\begin{equation}\label{trans:high-ranks}
 \dim B_a^j
 \le C\sum_i\#\{m\ge1:m\ell_i<a,\ |j-m\Delta_i|\le C_{\rm loc}+n\}.
\end{equation}
This follows by summing the crossing exact-sequence bounds, starting
with the zero relative group below the first event.
If $\Delta_i>0$, the number in braces is bounded by
$1+2(C_{\rm loc}+n)/\Delta_i$, up to an integer rounding constant.
If $\Delta_i=0$, it is zero once $j>C_{\rm loc}+n$.
The finite prime list therefore gives a uniform bound on
$\dim B_a^j$ in every sufficiently high degree.

The fixed low-degree interval needs a different argument, which also
accounts for all zero-mean-index iterates.  Choose a regular
$a_+\in(a+D,a+D+1)$.  By ordinary death,
\begin{equation}\label{trans:kernel-window}
 \ker\bigl(H_j(X_a,M)\to H_j(\Lambda,M)\bigr)
 =\ker\bigl(H_j(X_a,M)\to H_j(X_{a_+},M)\bigr).
\end{equation}
The kernel on the right is bounded in dimension by the sum of local
ordinary ranks in this interval of cut values of width at most $D+1$.  The number
of its events is at most
\[
 \sum_i\left(2+\frac{D+1}{\ell_i}\right),
\]
independently of $a$.  Uniform local ranks bound the kernel, and the
global-image rank is bounded by the fixed global degree rank.
Consequently $\dim H^j(X_a,M)$ is bounded uniformly in $a$ for each
fixed low degree $j$.

The sign-line Gysin sequence gives
\[
 \dim B_a^j\le\dim B_a^{j-1}+\dim H^j(X_a,M),
 \qquad B_a^{-1}=0.
\]
Induction, including degree zero, bounds $B_a^j$ uniformly throughout
any fixed low-degree interval.  Combining this with
\eqref{trans:high-ranks}, every degree in the finitely many remaining
bands has uniformly bounded dimension.  Since the number of such
degrees is also bounded, their sum bounds $T(a)$ as asserted.

Cuts $a$ in a bounded interval have only finitely many crossing stages,
so they satisfy a common bound as well.  Finally, if $a$ is regular but
$2a$ is not, choose a nearby $a'$ in the same no-critical interval such
that $2a'$ is regular.  The critical values and their halves are locally
finite.  The no-critical restriction isomorphism between the root
sublevels is natural with respect to global restriction, and hence
preserves both $B_a$ and its global image.  In particular
$T(a)=T(a')$.  This proves \eqref{trans:total-bound} for every positive
regular cut.
\end{proof}

We have bounded the total excess dimension at each regular cut without
assuming monotonicity as the cut varies.  The next section proves
monotonicity under doubling.  Together, these statements will force an
interval of maximal transient dimension on which the activation of a
new global class is incompatible with the local evaluation-cup relation.
